\section{Ordinary-line risk under the point law}\label{sec:riskproof}

Fix a root \(p\) and its deterministic deadline \(\tau=\tau_p\).
All observations below are stopped and saved at the root's first
positive boundary or global stop. A quantity explicitly restricted
to the active root is zero afterwards. Write
\(F_q(t)=\int_0^t m_u(q)\dd u\), and let
\[
 I_p=\int_0^{\theta_p\wedge T}\ind_{p\text{ active}}H_p(t)\dd t
\]
be its integrated ordinary-line hazard. Earlier root death has
\(\Prob^{(p)}\)-probability zero. We will bound \(I_p\) by a proxy
having sharp fixed-order moments, plus an error small in \(L^1\).
The distinction between those two kinds of estimate is essential.

\subsection{Removing ordinary-line credit at additive cost}
For an ordinary line through \(p\), let
\(H_\ell(t)=r_\ell(t)\sum_{q\in\ell\setminus\{p\},\ q\text{ active}}w_q(t)\).
The two-root, no-witness case of Proposition~\ref{prop:mixed}, followed
by the conditional density identity, gives
\begin{equation}\label{eq:conditionalH}
 \E^{(p)}[\ind_{p\text{ before boundary at }t}H_\ell(t)\mid\F_u]
 \le(1+o(1))\ind_{p\text{ before boundary at }u}H_\ell(u)
 \quad(u\le t).
\end{equation}
The correction ratio is bounded by a fixed power of
\((1-\epsilon_n)^{-1}\). Conditional division is performed only
where \(Z_p(u)>0\); its complement is null under the point law.
The auxiliary-root sum excludes \(p\), so no self-square is
silently included in this two-root argument.

One selected witness \(z\in\ell\) and two leaves give
\begin{equation}\label{eq:secondH}
 \E^{(p)}[\ind_{p\text{ before boundary},\ z\in S_t}H_\ell(t)^2]
 \le C L_\ell(L_\ell+B)t/n.
\end{equation}
For distinct leaves the connection factor is exactly \(r_\ell^2\);
a repeated leaf is paid using \(r^2w^2\le Brw\).
The one-leaf version is at most \(CL_\ell t/n\).
All labels here lie on the same ordinary line, so there is no
designated-line source multiplier.

Since \(r\) starts moving at the first line birth,
\(1-r_\ell(t)=\int_{\text{activation}}^t r_\ell\Lambda_\ell\dd u\).
The correct self-rate bound is
\begin{equation}\label{eq:selfrate}
 r_\ell\Lambda_\ell\le H_\ell/n+B/n.
\end{equation}
Multiplying the leaf sums in the hazard by \(r_\ell\) therefore loses
a nonnegative amount whose expectation is at most
\begin{equation}\label{eq:removecredit}
 Cn^{-3}\sum_{\ell\ni p}(L_\ell^3+BL_\ell^2)
 +CBn^{-3}\sum_{\ell\ni p}L_\ell^2
 =O(1+(B/n)\log n)=o(\log n).
\end{equation}
For completeness, insert the integral for \(1-r\), use
\eqref{eq:conditionalH} to move the future \(H\) to the earlier
time, cover activation by \(\sum_{z\in\ell}\ind_{z\in S_u}\),
and apply \eqref{eq:secondH}. The second term in
\eqref{eq:removecredit} is the separate self-rate contribution from
\eqref{eq:selfrate}. The line sums are Lemma~\ref{lem:lines}.
This is an additive error, not a constant-factor loss in the leading
\(\log n\) coefficient.

\subsection{Long lines: replace the leaf sum, retain the birth time}
Call a line long when \(L_\ell\ge B\log n\). Fix a selected witness
\(z\in\ell\), and put
\[
 H=r_\ell\sum_{q\in\ell\setminus\{p\},\ q\text{ active}}m_t(q)w_q,
 \qquad D_z(t)=\sum_{q\in\ell\setminus\{p,z\}}m_t(q).
\]
For every fixed integer \(k\), the mixed hierarchy gives
\begin{equation}\label{eq:longmoment}
 \E^{(p)}[\ind_{p\text{ before boundary},\ z\in S_t}H^k]
 \le(1+o(1))\frac{F_z(t)}n
       [D_z(t)^k+C_kB L_\ell^{k-1}].
\end{equation}
Indeed, the all-distinct leaves have exactly the required connection
\(r_\ell^k\) and unit source price. A term with \(a<k\) distinct
leaves costs \(B^{k-a}L_\ell^a\le BL_\ell^{k-1}\), since
\(B\le L_\ell\). There are finitely many repetition patterns for
fixed \(k\). No division by a positive lower bound on \(F_z\) is
used; if it is zero the corresponding birth event is null.

For fixed \(\epsilon>0\), use
\(H\ind_{H>D+\epsilon L}\le H^k/(D+\epsilon L)^{k-1}\).
Since \(D\le L\), the main term is at most
\(L(1+\epsilon)^{-(k-1)}\), while the repeated-leaf term is
\(C_{k,\epsilon}B\). Sum over witnesses, long lines and time,
with the hazard factor \(1/n\). As \(B/L\le1/\log n\), the
positive error \(R_{\rm long}\) made by replacing the leaf sum by
\(D_z(t)\) satisfies
\begin{equation}\label{eq:longerror}
 \limsup_{n\to\infty}\sup_p
 \frac{\E^{(p)}R_{\rm long}}{\log n}
 \le C\epsilon+C(1+\epsilon)^{-(k-1)}.
\end{equation}
First choose \(\epsilon\), next a fixed sufficiently large \(k\),
and only then let \(n\) grow. This proves an \(o(\log n)\)
\(L^1\) error for the same process.

For a witness born at time \(u\), define its deterministic remaining
cost and birth price by
\[
 A_z(u)=\frac1n\int_u^\tau D_z(t)\dd t\le\frac{TL_{pz}}n,
 \qquad b_z=\frac1n\int_0^\tau m_u(z)A_z(u)\dd u.
\]
Let \(P_{\rm long}\) be the sum of \(A_z\) over long-line witnesses
born before the root boundary, and let \(M_{\rm long}=\sum_z b_z\).
Insert \(A_z(u)/(TL_{pz}/n)\le1\) as a deterministic birth-retention
factor in the mixed hierarchy. It follows that for every fixed \(m\),
\begin{equation}\label{eq:longproxy}
 \|P_{\rm long}\|_{L^m(\Prob^{(p)})}
 \le M_{\rm long}+o(\log n).
\end{equation}
Here are the non-generic terms in this assertion. A repeated witness
block of size \(r\ge2\) has total price at most
\(C_r n^{-(r+1)}\sum_\ell L_\ell^{r+1}=O_r(1)\), whereas a
singleton block costs \(O(\log n)\). Thus repetitions lose a power
of \(\log n\). For two distinct witnesses sharing a designated
line, that line cannot pass through the root. By
\eqref{eq:fixedline}, the weighted configuration count is
\[
 \sum_{z_1}L_{pz_1}
 \sum_{\substack{z_2\text{ on a designated}\\\text{line through }z_1}}L_{pz_2}
 \le C n^{3+\xi}\log n.
\]
The two birth-and-cost factors supply \(n^{-4}\), giving
\(O(n^{-1+\xi}\log n)\). Remaining singleton factors cost
\(O(\log n)\) each. A finite union over pairs covers multiple
exceptions. The all-distinct, unshared case keeps the exact product
of the \(b_z\), with leading factor \(1+o(1)\).

\subsection{Short stars and their exceptional configurations}
An active short line through \(p\) has one selected witness \(z\).
Each active leaf \(q\ne p\) on that line gives a star \((p,q,z)\).
Let
\[
 X_t=\sum_{\text{short stars}}m_t(q)r_{\ell(p,q)}w_q,
 \qquad
 D_t=\frac1n\sum_{\text{deterministic short stars}}m_t(q)F_z(t).
\]
The actual sum is zero after the root boundary. In the deterministic
sum, \(p,q,z\) are distinct and collinear on a short ordinary line;
there is no selected-point condition.
Put \(t_0=n^{\gamma-1}\).

\begin{lemma}[Short-star moment estimate]\label{lem:short}
For each fixed positive integer \(m\), uniformly for \(t_0\le t\le T\)
and all roots,
\begin{equation}\label{eq:shortmoment}
 \E^{(p)}X_t^m\le(1+o(1))D_t^m+\epsilon_m(n)(tn\log n)^m,
 \qquad \epsilon_m(n)\longrightarrow0.
\end{equation}
\end{lemma}
\begin{proof}
In an \(m\)-fold product let \(a\) be the number of distinct leaves
and \(j\) the number of distinct witnesses. The active root makes
the selected witness on each of its lines unique, so
\(1\le j\le a\le m\). A leaf coinciding with a witness or the root
gives zero. If \(j<m\), the first leaf and witness of each star
line have \(O(n^2\log n)\) choices by \eqref{eq:line2}.
Additional leaves on that line have at most \(B\log n\) choices,
and repeated weights cost \(B^{m-a}\). The source price is
\(C_m(t/n)^j\). Dividing this bound by \((tn\log n)^m\) gives
\[
 C_m\left(\frac B{nt}\right)^{m-j}(\log n)^{a-m}=o(1),
\]
uniformly since \(B/(nt)\le n^{d-\gamma}\).

It remains to treat \(j=a=m\), with distinct star lines. The
following list exhausts the extra connections among different stars.
\begin{enumerate}[label=(\roman*)]
\item Two labels share a designated line. Fix the first label.
The line is outside the root, so \eqref{eq:fixedline} bounds the
second label together with its star partner by \(C n^{1+\xi}\),
instead of the generic \(O(n^2\log n)\) row. This covers leaf--leaf,
leaf--witness and witness--witness sharing.
\item A line connecting two leaves has length greater than
\(n^\gamma\). For a fixed first leaf there are
\(O(n^{2-\gamma})\) possible second leaves: sum the height counts
up to \(O(n^{1-\gamma})\). Its star witness has at most
\(B\log n\) choices. Relative to a generic star row the saving
is \(O(n^{d-\gamma})\).
\item A leaf--leaf line of length at most \(n^\gamma\) is active
through a new witness. There are at most \(n^\gamma\) positions
for it and an additional source price \(C_mt/n\), giving
\(O(n^{\gamma-1})\).
\item That line is active through an existing third star witness.
It cannot be either leaf's own witness, since then their star lines
would coincide. For a third witness use
\(\sum_{z\in\ell(q_i,q_j)}L_{pz}\le C n^{1+\xi}\), keeping all
\(m\) original birth prices. This again saves \(n^{-1+\xi}\).
An active leaf or the root cannot itself be a selected witness.
\end{enumerate}
Outside these cases the additional ordinary-line credits are all
one, there are no designated-line connections, and the star credits
exactly match the mixed-moment factors. The source price is the
exact product of \(F_z(t)/n\); keeping the masks produces \(D_t^m\).
Each single exception is bounded after dropping other restrictions,
so a finite union over label pairs also covers overlapping exceptions.
Removing remaining factors costs only \(C_m\), including
\(\rho^{-O(m^2)}\). The savings above and the correction
\((1-\epsilon_n)^{-O(m^2)}=1+o(1)\) prove the assertion.
\end{proof}

Minkowski's inequality now yields
\begin{equation}\label{eq:shortintegral}
 \left\|\int_{t_0}^\tau X_t/n\dd t\right\|_m
 \le\int_{t_0}^\tau D_t/n\dd t+o(\log n),
\end{equation}
with the integral zero when \(\tau<t_0\).
For earlier times the one-witness, one-leaf bound gives
\(\E^{(p)}X_t=O(tn\log n)\). Its integrated contribution is
\(O(t_0^2\log n)=o(\log n)\) in \(L^1\). No higher moment for
this early-time error is asserted.

\subsection{Hazard concentration and the weight cap}
Collect the credit-removal error, long-line leaf replacement error,
and early short-star contribution into \(R_n\ge0\), so
\(\sup_p\E^{(p)}R_n=o(\log n)\). Define
\[
 Y_n=P_{\rm long}+\int_{t_0}^\tau X_t/n\dd t.
\]
Then \(I_p\le Y_n+R_n\), and \eqref{eq:longproxy} and
\eqref{eq:shortintegral} give, for every fixed \(m\),
\[
 \|Y_n\|_m\le M_{\rm long}+M_{\rm short}+o(\log n)
 \le\mathcal B_n^{\Gamma^c}(p;F_\tau)+o(\log n).
\]
The last inequality is \eqref{eq:timesym}: the division into long
and short lines is symmetric in leaf and witness, and all four
fields use the same \(\tau\).
Proposition~\ref{prop:pair} and \eqref{eq:riskinput} therefore give
\(\|Y_n\|_m\le(r_*+o(1))\log n\), uniformly in the root.

For any \(b>r_*\), choose \(\epsilon>0\) with
\(r_*+\epsilon<b-\epsilon\). Markov on \(R_n\), followed by
the fixed-\(m\) moment bound, shows
\[
 \limsup_{n\to\infty}\sup_p\Prob^{(p)}(I_p>b\log n)
 \le\left(\frac{r_*+\epsilon}{b-\epsilon}\right)^m.
\]
The right side tends to zero as the fixed proof order \(m\) is
chosen larger. This quantifier order does not require constants
uniform over all \(m\), or an order depending on \(n\).
By \eqref{eq:product}, hitting \(B=n^d\) requires
\(I_p\ge d\log n-O(1)\). Choose \(r_*<b<d\). We have proved
\begin{equation}\label{eq:captail}
 \sup_p\Prob^{(p)}(\text{weight cap is an early saving cause})=o(1).
\end{equation}

